% English version of cv_4.tex
\documentclass{resume}
\usepackage[left=0.45in,top=0.35in,right=0.45in,bottom=0.35in]{geometry}
\usepackage{fontspec}
\setmainfont{LinLibertine_R.otf}[
	BoldFont=LinLibertine_RB.otf,
	ItalicFont=LinLibertine_RI.otf,
	BoldItalicFont=LinLibertine_RBI.otf
]
\usepackage{enumitem}
\usepackage{hyperref}
\hypersetup{colorlinks=false,pdfborder={0 0 0},pdftitle={Peihuan Ni - Curriculum Vitae},pdfauthor={Peihuan Ni}}
\setlist[itemize]{leftmargin=1.3em,itemsep=1pt,topsep=2pt,parsep=0pt,partopsep=0pt}
\newenvironment{publicationentry}{\begin{list}{}{\setlength{\leftmargin}{0.6em}\setlength{\rightmargin}{0em}\setlength{\topsep}{3pt}\setlength{\partopsep}{0pt}\setlength{\parsep}{0pt}\setlength{\itemsep}{0pt}}\item[]}{\end{list}}
\makeatletter
\renewcommand{\printname}{\begingroup\centering{\MakeUppercase{\namesize\bfseries \@name}}\par\endgroup\nameskip}
\makeatother
\name{Peihuan Ni}
\address{Institute of Automation, Chinese Academy of Sciences}
\address{95 Zhongguancun East Road, Haidian District, Beijing, China}
\address{+86 158 0375 7631 \\ \href{mailto:nipeihuan24@mails.ucas.ac.cn}{nipeihuan24@mails.ucas.ac.cn} \\ \href{https://peihuanni.github.io}{peihuanni.github.io}}
\setlength{\parskip}{4pt}
\renewcommand{\sectionskip}{\vspace{8pt}}
\renewcommand{\sectionlineskip}{\vspace{5pt}}
\begin{document}
\fontsize{10}{12}\selectfont
\begin{rSection}{Research Interests}
AI infrastructure and distributed systems for large-model training and inference; efficient inference architectures for large language models and vision-language-action models; hardware accelerators and embodied artificial intelligence.
\end{rSection}

\begin{rSection}{Ph.D. Applications}
Seeking Ph.D. opportunities in China and abroad for 2027 entry, with research interests in ML systems for training and inference acceleration, AI infrastructure, and computer architecture.
\end{rSection}

\begin{rSection}{Education}
\textbf{University of Chinese Academy of Sciences} \hfill Sep. 2024 -- Jun. 2027 (expected)\\
Master's student in Artificial Intelligence, School of Artificial Intelligence\\
Institute of Automation, Chinese Academy of Sciences; Efficient Intelligent Computing and Learning Lab (CLab)\\
Advisors: Jian Cheng and Gang Li \hfill GPA: 3.81

\textbf{Jilin University} \hfill Sep. 2020 -- Jun. 2024\\
Bachelor's degree in Microelectronics Science and Engineering\\
College of Electronic Science and Engineering \hfill Advisor: Kaige Yan\\
College English Test: CET-4: 603; CET-6: 553
\end{rSection}

\begin{rSection}{Research Experience}
\textbf{First-Author Publications}
\begin{publicationentry}
\textbf{Peihuan Ni}, Zitao Mo, Tielong Liu, Hongli Wen, Zeyu Zhu, Minnan Pei, Junwen Si, Weifan Guan, Peisong Wang, Qinghao Hu, Gang Li and Jian Cheng.
\textbf{APEX: Integer-only Non-linear Function Approximation for Efficient Cross-Modal Inference.}
In \textbf{DATE 2026}. (CCF-B)
\begin{itemize}
\item Identified substantial differences in quantized data distributions across model modalities, leading to varying sensitivity to nonlinear approximation errors; retaining high precision uniformly incurs excessive computational bit widths and hardware costs.
\item Proposed a unified integer-only architecture that merges the computation graphs of common nonlinear operators and uses bit-level pruning to reduce bit widths while preserving model accuracy.
\item Improved task accuracy by up to approximately 0.7\% for language models and 1.3\% for vision models, while enabling near-lossless inference for multimodal models.
\item Achieved area reduction factors of 1.73--8.71$\times$ and power reduction factors of 1.21--10.83$\times$ relative to existing methods.
\end{itemize}
\end{publicationentry}
\begin{publicationentry}
\textbf{Peihuan Ni}, Zeyu Zhu, Weifan Guan, Gang Li, Zitao Mo, Tielong Liu, Junwen Si, Qinghao Hu, Peisong Wang and Jian Cheng.
\textbf{ACE-VLA: Action-Centric Evaluation for Token Pruning in Vision-Language-Action Models.}
In \textbf{CoRL 2026}. (BAAI Class A)
\begin{itemize}
\item Identified that existing VLA pruning methods largely adopt text--vision attention metrics from VLMs, underusing task-phase and motion semantics encoded in action outputs.
\item Proposed an action-centric gradient-based importance metric and a motion-trajectory-based dynamic token pool to reduce importance evaluation overhead.
\item Modeled robot motion trends from a kinematic perspective to dynamically adjust pruning ratios beyond instantaneous motion states.
\item Improved task success rates by approximately 1.4\% on $\pi_{0.5}$ and 2.2\% on $\pi_0$, with 1.33--1.57$\times$ inference speedups; outperformed existing baselines in overall performance on OpenVLA-OFT.
\item Validated transferability through real-robot experiments and improved the interpretability of VLA token pruning.
\end{itemize}
\end{publicationentry}

% \newpage

\textbf{Co-Authored Publications}
\begin{publicationentry}
Zeyu Zhu, Gang Li, Minnan Pei, Zitao Mo, \textbf{Peihuan Ni}, Peisong Wang, Tielong Liu and Jian Cheng.
\textbf{KL-MoE: A Hierarchical MoE Pruning Framework Exploiting KL Divergence.}
In \textbf{DAC 2026}. (CCF-A)
\end{publicationentry}
\begin{publicationentry}
Tielong Liu, Zeyu Zhu, Gang Li, Zitao Mo, Chen Bai, Junwen Si, \textbf{Peihuan Ni}, Minnan Pei, Peisong Wang, Zhuoran Song, Xiaoyao Liang and Jian Cheng.
\textbf{FLUTE: A LUT-Based Floating-Point Engine for Efficient LLM Acceleration Exploiting Core-Tail Hierarchical Alignment.}
In \textbf{MICRO 2026}. (CCF-A)
\end{publicationentry}
\begin{publicationentry}
Junwen Si, Gang Li, Tielong Liu, \textbf{Peihuan Ni}, Minnan Pei, Zeyu Zhu, Zitao Mo, Peisong Wang, Zhongfeng Wang and Jian Cheng.
\textbf{DOMI: A Dataflow Optimization Framework for Efficient Mamba Inference on FPGAs.}
Submitted to \textbf{TCAD}. (CCF-A)
\end{publicationentry}
\begin{publicationentry}
Zeyu Zhu, Gang Li, Peisong Wang, \textbf{Peihuan Ni}, Tielong Liu, Minnan Pei, Junwen Si, Zhuoran Song, Xiaoyao Liang and Jian Cheng.
\textbf{DALI: A Workload-Aware Offloading Framework for Efficient MoE Inference on Local PCs.}
Submitted to \textbf{ATC 2026}. (CCF-A)
\end{publicationentry}
\end{rSection}

\begin{rSection}{Internship Experience}
\textbf{Huawei, ICT BG} \hfill Jul. 2026 -- Jul. 2027\\
\textit{LLM Inference System Acceleration on the Ascend 950 NPU}
\begin{itemize}
\item Analyze and optimize LLM inference on the Ascend 950 NPU, focusing on operator scheduling, on-chip dataflow, and compute resource utilization.
\item Investigate bottlenecks across Cube and Vector units and the memory hierarchy, and explore system-level optimizations for MoE and long-context inference.
\end{itemize}
\end{rSection}

\begin{rSection}{Projects}
\textbf{VLA Model Quantization and FPGA Edge Deployment} \hfill \textbf{Project Lead}
\begin{itemize}
\item Designed a W4A8 quantization scheme with 4-bit weights and 8-bit activations and KV cache, combined with integer-only nonlinear approximations for near-lossless edge inference.
\item Implemented fused quantization and dequantization operators using FP16--INT8 multiply-accumulate (MAC) units, reducing intermediate writebacks and off-chip memory accesses while keeping more activations on chip.
\item Reused FP16--INT8 MAC units for operators such as RoPE and reordered selected computations to improve compute-unit and on-chip data reuse.
\item Designed bit-serial compute units for W4A8 mixed-precision matrix multiplication, a KV reordering module supporting different attention head dimensions, and crossbar-based data mapping from BRAM to the compute array.
\item Designed a unified integer-only nonlinear unit supporting multiple operators with reduced hardware area and power consumption; results published at DATE 2026.
\item Contributed to VLIW compiler design, with responsibility for on-chip dataflow, the instruction processor, and operator control modules; reduced the instruction count per model inference by over 32$\times$ using address offsets and multiple offset modes.
\item Built a GPU-based bit-accurate simulator and verified consistency with Vivado module-level simulations and full-model execution on FPGA.
\item Completed end-to-end FPGA deployment, including DMA memory-access debugging, HBM bandwidth optimization, VHK158 board data-path debugging, timing closure, and PetaLinux system builds.
\end{itemize}
\vspace{3pt}
\textbf{LLM Analysis and Inference Optimization on the Ascend A5 NPU} \hfill \textbf{Project Lead}
\begin{itemize}
\item Analyzed bottlenecks caused by asymmetric throughput between matrix and vector compute units in the Ascend architecture, identifying nonlinear operator overhead, data movement, and idle compute resources in MoE inference.
\item Explored fine-grained pipelining and overlapping of matrix multiplication and nonlinear computation to balance Cube and Vector workloads and improve overall hardware utilization.
\end{itemize}
\end{rSection}

\begin{rSection}{Honors and Awards}
Outstanding Student, University of Chinese Academy of Sciences (2025)\\
Third Prize, Intel Cup National Undergraduate Electronic Design Contest -- Invitational (2022)\\
Second Prize, China Undergraduate Mathematical Contest in Modeling (2022)\\
Third Prize, Chinese Mathematics Competitions for College Students (2023)
\end{rSection}
\begin{rSection}{Academic Service}
Reviewer for WACV, EMNLP, CoRL and HPCA.
\end{rSection}
\end{document}
